\documentclass[a4paper,11pt]{article}
\usepackage[utf8]{inputenc}
\usepackage[T1]{fontenc}
\usepackage[ngerman]{babel}
\usepackage{geometry}
\geometry{left=2.5cm, right=2.5cm, top=2.5cm, bottom=2.5cm}
\usepackage{longtable}
\usepackage{booktabs}
\usepackage{array}
\usepackage{xcolor}
\usepackage{colortbl}
\usepackage{hyperref}
\usepackage{fancyhdr}
\usepackage{tabularx}
\usepackage{enumitem}
\usepackage{graphicx}
\usepackage{listings}

\definecolor{headerblue}{HTML}{1565C0}
\definecolor{tableheader}{HTML}{E3F2FD}
\definecolor{pass}{HTML}{2E7D32}
\definecolor{warn}{HTML}{F57F17}
\definecolor{codebg}{HTML}{F5F5F5}

\hypersetup{
    colorlinks=true,
    linkcolor=headerblue,
    urlcolor=headerblue
}

\pagestyle{fancy}
\fancyhf{}
\fancyhead[L]{\textbf{INSPECTAIR}}
\fancyhead[R]{Coding Standards \& Rules v2.0}
\fancyfoot[C]{\thepage}
\renewcommand{\headrulewidth}{0.4pt}

\setlength{\parindent}{0pt}
\setlength{\parskip}{6pt}

\lstset{
    backgroundcolor=\color{codebg},
    basicstyle=\ttfamily\small,
    breaklines=true,
    frame=single,
    rulecolor=\color{gray!40},
    xleftmargin=2em,
    framexleftmargin=1.5em,
    language=C++,
    showstringspaces=false,
    tabsize=4,
    literate={ä}{{\"a}}1 {ö}{{\"o}}1 {ü}{{\"u}}1 {Ä}{{\"A}}1 {Ö}{{\"O}}1 {Ü}{{\"U}}1 {ß}{{\ss}}1
}

\begin{document}

\begin{center}
{\LARGE\textbf{INSPECTAIR}}\\[0.3cm]
{\Large Coding Standards \& Rules}\\[0.3cm]
{\large Luftqualitätsmessgerät mit ESP32-S3}\\[0.5cm]
\begin{tabular}{ll}
\textbf{Projekt:} & InspectAir -- Luftqualitätsmessgerät \\
\textbf{Version:} & 2.0 \\
\textbf{Datum:} & Februar 2026 \\
\textbf{Autoren:} & Maximilian Liebl, Cedric Stroh, Jannis Messer \\
\textbf{Modul:} & Mikrocomputertechnik 1, DHBW Ravensburg \\
\textbf{Dozent:} & Thomas Stingl (Airbus) \\
\end{tabular}
\end{center}

\vspace{0.5cm}
\noindent\textit{Dieses Dokument wurde unter Verwendung von KI-Tools (Claude, Anthropic) zur Überprüfung erstellt.}
\vspace{0.3cm}

\tableofcontents
\newpage

% ============================================================
\section{Allgemeine Regeln}
% ============================================================

\begin{itemize}[leftmargin=2em]
    \item \textbf{Sprache:} C++ (Arduino Framework via PlatformIO)
    \item \textbf{Einrückung:} 4 Spaces (keine Tabs)
    \item \textbf{Zeilenlänge:} max. 120 Zeichen
    \item \textbf{Encoding:} UTF-8
    \item Alle Dateien enden mit einer Leerzeile
\end{itemize}

% ============================================================
\section{Namenskonventionen}
% ============================================================

\begin{longtable}{p{3.5cm} p{4.5cm} p{5cm}}
\toprule
\rowcolor{tableheader}
\textbf{Element} & \textbf{Konvention} & \textbf{Beispiel} \\
\midrule
\endhead
Variablen        & \texttt{snake\_case}       & \texttt{sensor\_data} \\
Konstanten       & \texttt{UPPER\_SNAKE\_CASE} & \texttt{CO2\_THRESHOLD} \\
Funktionen       & \texttt{snake\_case}       & \texttt{sensors\_init()} \\
Structs/Enums    & \texttt{PascalCase}        & \texttt{SensorData} \\
Dateien          & \texttt{snake\_case.cpp/.h} & \texttt{sensors.cpp} \\
Defines          & \texttt{UPPER\_SNAKE\_CASE} & \texttt{\#define PIN\_SDA 8} \\
\bottomrule
\end{longtable}

% ============================================================
\section{Dokumentation (Doxygen)}
% ============================================================

Alle öffentlichen Funktionen und Structs müssen dokumentiert werden:

\begin{lstlisting}
/**
 * @brief Initialisiert alle Sensoren
 *
 * Initialisiert I2C, UART und alle angeschlossenen Sensoren.
 * Bei Fehler wird das entsprechende ok-Flag auf false gesetzt.
 *
 * @param config Zeiger auf Konfigurationsstruktur
 * @return true wenn mindestens ein Sensor verfuegbar
 * @return false wenn kein Sensor initialisiert werden konnte
 */
bool sensors_init(const SensorConfig* config);
\end{lstlisting}

% ============================================================
\section{Fehlerbehandlung}
% ============================================================

\begin{itemize}[leftmargin=2em]
    \item Jeder Sensor hat ein eigenes \texttt{ok}-Flag im \texttt{SensorData} Struct
    \item Sensor-Fehler werden geloggt aber nicht als fatal behandelt
    \item Display zeigt \texttt{-{}-} für nicht verfügbare Werte
    \item Kritische Fehler (Display-Init) führen zu Halt
\end{itemize}

% ============================================================
\section{Git Workflow}
% ============================================================

\begin{itemize}[leftmargin=2em]
    \item \textbf{Main-Branch:} Nur getesteter, funktionierender Code
    \item \textbf{Feature-Branches:} \texttt{feature/sensor-module}, \texttt{feature/ui-detailed}
\end{itemize}

Commit-Messages mit Präfix:

\begin{longtable}{p{3cm} p{10cm}}
\toprule
\rowcolor{tableheader}
\textbf{Präfix} & \textbf{Bedeutung} \\
\midrule
\endhead
\texttt{feat:}     & Neue Funktionalität \\
\texttt{fix:}      & Bugfix \\
\texttt{docs:}     & Dokumentation \\
\texttt{refactor:} & Code-Umstrukturierung \\
\bottomrule
\end{longtable}

% ============================================================
\section{Verbotene Praktiken}
% ============================================================

\begin{itemize}[leftmargin=2em]
    \item[\textcolor{red}{\textbf{X}}] Keine Magic Numbers -- immer \texttt{\#define} oder \texttt{const} verwenden
    \item[\textcolor{red}{\textbf{X}}] Keine globalen Variablen ohne \texttt{static} (Datei-Scope begrenzen)
    \item[\textcolor{red}{\textbf{X}}] Keine \texttt{delay()} in \texttt{loop()} -- \texttt{millis()} oder \texttt{yield()} verwenden
    \item[\textcolor{red}{\textbf{X}}] Keine \texttt{String}-Klasse -- \texttt{char[]} Buffer verwenden
    \item[\textcolor{red}{\textbf{X}}] Keine dynamische Speicherallokation in \texttt{loop()}
    \item[\textcolor{red}{\textbf{X}}] Kein \texttt{Serial.print()} in ISRs
\end{itemize}

% ============================================================
\section{ESP32/FreeRTOS-spezifische Regeln}
% ============================================================

\begin{itemize}[leftmargin=2em]
    \item \texttt{volatile} für alle Variablen die in ISRs gelesen/geschrieben werden
    \item \texttt{IRAM\_ATTR} für Interrupt Service Routines
    \item Task-Stack-Größe beachten (Standard: 4KB, bei viel lokalem Speicher erhöhen)
    \item \texttt{portENTER\_CRITICAL()} / \texttt{portEXIT\_CRITICAL()} für shared resources zwischen Tasks
\end{itemize}

% ============================================================
\section{Const Correctness}
% ============================================================

\begin{itemize}[leftmargin=2em]
    \item Pointer-Parameter die nicht geändert werden: \texttt{const char* text}
    \item Methoden die Objekt nicht ändern: \texttt{int getValue() const}
    \item Lokale Variablen die nicht geändert werden: \texttt{const int max\_retries = 10}
\end{itemize}

% ============================================================
\section{Defensive Programmierung}
% ============================================================

\begin{itemize}[leftmargin=2em]
    \item Null-Pointer-Checks vor Dereferenzierung
    \item Array-Bounds validieren
    \item Timeout bei blockierenden Operationen (WiFi, Sensor-Reads)
    \item Sinnvolle Default-Werte bei fehlgeschlagenen Reads
\end{itemize}

% ============================================================
\section{Speicher-Management}
% ============================================================

\begin{itemize}[leftmargin=2em]
    \item Statische Allokation bevorzugen
    \item Dynamische Allokation nur in \texttt{setup()} / \texttt{begin()} Funktionen
    \item Buffer-Größen als Konstanten definieren
    \item Bei Arrays: \texttt{sizeof()} statt hardcoded Größen
\end{itemize}

\end{document}
